\documentclass[11pt,a4paper]{article}

% --- ARXIV PREPRINT PACKAGES ---
\usepackage[utf8]{inputenc}
\usepackage[T1]{fontenc}
\usepackage{lmodern}
\usepackage[indonesian]{babel}
\usepackage[margin=1in, top=1.1in, bottom=1.1in, headheight=25pt]{geometry}
\usepackage{amsmath, amsfonts, amssymb, bm}
\usepackage{graphicx}
\usepackage{booktabs}
\usepackage{tabularx}
\usepackage{longtable}
\usepackage{array}
\usepackage{multirow}
\usepackage{xcolor}
\usepackage{hyperref}
\usepackage{caption}
\usepackage{subcaption}
\usepackage{float}
\usepackage{fancyhdr}
\usepackage[nopatch=footnote]{microtype}
\usepackage{authblk}
\usepackage{enumitem}

% --- COLOR DEFINITIONS ---
\definecolor{arxivblue}{RGB}{0, 51, 153}
\definecolor{headergray}{RGB}{90, 100, 110}
\definecolor{darkslate}{RGB}{30, 41, 59}

% --- HYPERLINK SETUP ---
\hypersetup{
    colorlinks=true,
    linkcolor=arxivblue,
    citecolor=arxivblue,
    urlcolor=arxivblue,
    pdftitle={Analisis Curah Hujan Multi-Skala: Satelit Kebumen vs AWS IoT Jerukagung},
    pdfauthor={Tim Peneliti Presipitasi & AI Kebumen}
}

% --- GRAPHICS PATH ---
\graphicspath{
    {./Hasil_Analisis/}
    {Hasil_Analisis/}
    {./}
}

% --- ARXIV PREPRINT HEADER & FOOTER ---
\setlength{\headheight}{25pt}
\pagestyle{fancy}
\fancyhf{}
\fancyhead[L]{\small\textsf{\color{headergray}\textbf{A PREPRINT} --- ANALISIS CURAH HUJAN: PER JAM VS PER HARI}}
\fancyhead[R]{\small\textsf{\color{headergray}\thepage}}
\fancyfoot[C]{\footnotesize\textsf{\color{headergray}Laboratorium Sains Atmosfer \& AI Geospasial Kebumen}}
\renewcommand{\headrulewidth}{0.4pt}

% --- SECTION STYLING ---
\usepackage{titlesec}
\titleformat{\section}{\large\bfseries\color{arxivblue}}{\thesection.}{0.5em}{}
\titleformat{\subsection}{\normalsize\bfseries\color{darkslate}}{\thesubsection}{0.5em}{}

% --- TITLE & AUTHOR ---
\title{\vspace{-1.2cm}\textbf{\Large Analisis Curah Hujan Multi-Skala: Komparasi 8 Produk Presipitasi Satelit \& Reanalisis Kebumen terhadap Pengamatan AWS IoT Jerukagung pada Resolusi Per Jam (\textit{Hourly}) dan Per Hari (\textit{Daily})}}

\author[1]{\textbf{Tim Peneliti Presipitasi \& AI Kebumen}\thanks{Email korespondensi: \texttt{penelitian.ai-cuaca@kebumen-project.org}}}
\affil[1]{\small Laboratorium Kecerdasan Buatan Terapan \& Sains Atmosfer Geospasial, Kebumen}

\date{\small\today}

\begin{document}

\maketitle

\begin{abstract}
\noindent Validasi dan karakterisasi estimasi presipitasi satelit serta model reanalisis atmosfer terhadap observasi stasiun darat otomatis (\textit{Automatic Weather Station} / AWS IoT) memerlukan pemahaman mendalam pada skala temporal yang berbeda. Laporan ini menyajikan analisis komparasi multi-skala antara dataset curah hujan harian Kabupaten Kebumen (\texttt{Data\_Curah\_Hujan\_Kebumen.csv} yang mencakup 8 produk: \texttt{CHIRPS\_RNL}, \texttt{CHIRPS\_SAT}, \texttt{CHIRPS\_FNL}, \texttt{GSMaP}, \texttt{IMERG}, \texttt{PERSIANN}, \texttt{ERA5}, dan \texttt{ERA5\_LAND}) serta dataset jam-jaman terhadap stasiun AWS IoT Jerukagung sepanjang 563 hari valid sinkron (13.512 jam pengamatan). Hasil evaluasi mengungkap lonjakan akurasi yang sangat signifikan: pada skala per jam, presipitasi satelit memiliki korelasi moderat ($r = 0.414$ untuk IMERG dan $r = 0.313$ untuk GSMaP) akibat adanya \textit{spatial mismatch} dan \textit{sub-hourly lag}; namun pada skala harian, korelasi presipitasi satelit melonjak drastis hingga $r = \mathbf{0.619}$ ($\rho = \mathbf{0.642}$) untuk NASA GPM IMERG dan $r = \mathbf{0.513}$ ($\rho = \mathbf{0.595}$) untuk JAXA GSMaP. Analisis siklus diurnal 24-jam membuktikan bahwa puncak hujan konvektif di Kebumen terkonsentrasi pada sore hari pukul 15:00--18:00 WIB, yang berhasil ditangkap dengan sangat baik oleh satelit NASA IMERG dan AWS IoT.
\end{abstract}

\vspace{0.2cm}
\noindent\textbf{\textit{Keywords:}} Analisis Curah Hujan, AWS IoT Jerukagung, Data Curah Hujan Kebumen, NASA GPM IMERG, JAXA GSMaP, CHIRPS, ECMWF ERA5, Evaluasi Multi-Skala.

\vspace{0.5cm}
\hrule
\vspace{0.5cm}

% =============================================================================
\section{Pendahuluan \& Metodologi Dataset}
% =============================================================================
Analisis presipitasi dilakukan pada dua domain waktu:
\begin{enumerate}[leftmargin=*]
    \item \textbf{Skala Resolusi Per Jam (\textit{Hourly})}: Menggunakan 13.512 jam observasi sinkron 1-jam antara AWS IoT, JAXA GSMaP, NASA IMERG, dan ERA5 Hourly.
    \item \textbf{Skala Resolusi Per Hari (\textit{Daily})}: Mengintegrasikan dataset harian 8 produk satelit Kebumen (\texttt{Data\_Curah\_Hujan\_Kebumen.csv}) dengan total hujan harian stasiun AWS IoT Jerukagung sepanjang 563 hari valid.
\end{enumerate}

\begin{table}[htbp]
\centering
\small
\caption{Evaluasi Akurasi 8 Produk Presipitasi Harian Kebumen vs Stasiun AWS IoT Harian}
\label{tab:daily_eval}
\begin{tabular*}{\textwidth}{@{\extracolsep{\fill}}lrrrrrr@{}}
\toprule
\textbf{Produk Presipitasi} & \textbf{Pearson $r$} & \textbf{Spearman $\rho$} & \textbf{RMSE (mm/hari)} & \textbf{MAE (mm/hari)} & \textbf{PBIAS (\%)} & \textbf{KGE} \\
\midrule
\texttt{NASA GPM IMERG} & \textbf{0.619} & \textbf{0.642} & \textbf{13.57} & \textbf{7.64} & +8.5\% & \textbf{0.548} \\
\texttt{JAXA GSMaP}     & 0.513 & 0.595 & 18.34 & 8.22 & -6.2\% & 0.482 \\
\texttt{CHIRPS\_SAT}    & 0.498 & 0.581 & 17.89 & 8.05 & -4.8\% & 0.465 \\
\texttt{CHIRPS\_RNL}    & 0.485 & 0.570 & 18.12 & 8.19 & -5.1\% & 0.451 \\
\texttt{CHIRPS\_FNL}    & 0.472 & 0.558 & 19.04 & 8.52 & +12.4\% & 0.412 \\
\texttt{PERSIANN}       & 0.421 & 0.512 & 20.15 & 9.10 & -15.8\% & 0.380 \\
\texttt{ECMWF ERA5}     & 0.385 & 0.539 & 16.53 & 8.49 & +14.2\% & 0.365 \\
\texttt{ERA5\_LAND}     & 0.378 & 0.531 & 16.88 & 8.61 & +15.0\% & 0.354 \\
\bottomrule
\end{tabular*}
\end{table}

\begin{figure}[htbp]
    \centering
    \includegraphics[width=0.95\textwidth]{01_bar_evaluasi_8satelit_vs_aws_harian.png}
    \caption{Peringkat Koefisien Korelasi Pearson ($r$) dan Spearman ($\rho$) 8 Produk Satelit/Reanalisis Harian vs AWS IoT Jerukagung.}
    \label{fig:bar_daily}
\end{figure}

% =============================================================================
\section{Evaluasi Diagram Pencar Hexbin 8 Produk Presipitasi Harian}
% =============================================================================
Gambar \ref{fig:scatter_8sat} memperlihatkan sebaran data harian 8 produk satelit/reanalisis terhadap AWS IoT.

\begin{figure}[htbp]
    \centering
    \includegraphics[width=0.98\textwidth]{02_scatter_hexbin_8satelit_vs_aws_harian.png}
    \caption{Diagram Pencar Hexbin Density 8 Produk Presipitasi Satelit \& Reanalisis Kebumen vs AWS IoT Harian.}
    \label{fig:scatter_8sat}
\end{figure}

% =============================================================================
\section{Komparasi Multi-Skala: Resolusi Per Jam vs Resolusi Per Hari}
% =============================================================================
Gambar \ref{fig:comp_scale} dan \ref{fig:bar_jump} memperlihatkan perbandingan performa akurasi saat data dievaluasi pada skala 1-jam vs ketika diagregasikan ke skala 1-hari.

\begin{figure}[htbp]
    \centering
    \includegraphics[width=0.95\textwidth]{03_perbandingan_scatter_jam_vs_hari.png}
    \caption{Perbandingan Pencar Presipitasi: Skala 1-Jam (Atas) vs Skala 1-Hari (Bawah).}
    \label{fig:comp_scale}
\end{figure}

\begin{figure}[htbp]
    \centering
    \includegraphics[width=0.95\textwidth]{04_bar_lonjakan_akurasi_jam_vs_hari.png}
    \caption{Lonjakan Koefisien Korelasi Parametrik ($r$) dan Non-Parametrik ($\rho$) Saat Data Diagregasikan ke Skala Harian.}
    \label{fig:bar_jump}
\end{figure}

% =============================================================================
\section{Karakteristik Diurnal 24-Jam \& Skor Kontingensi Harian}
% =============================================================================
Gambar \ref{fig:diurnal} dan \ref{fig:contingency} menyajikan siklus diurnal 24-jam dan performa deteksi hujan berdasarkan ambang batas intensitas ($0.1 - 50.0\text{ mm/hari}$).

\begin{figure}[htbp]
    \centering
    \includegraphics[width=0.95\textwidth]{05_siklus_diurnal_24jam_cuaca_hujan.png}
    \caption{Siklus Diurnal 24-Jam Suhu, Kelembaban, Curah Hujan, dan Tekanan Permukaan di Stasiun Jerukagung Kebumen.}
    \label{fig:diurnal}
\end{figure}

\begin{figure}[htbp]
    \centering
    \includegraphics[width=0.90\textwidth]{06_skor_kontingensi_deteksi_hujan_harian.png}
    \caption{Kurva Metrik Kontingensi (CSI, POD, FAR, HSS) Deteksi Hujan Harian terhadap AWS IoT.}
    \label{fig:contingency}
\end{figure}

% =============================================================================
\section{Kurva Massa Ganda \& Konsistensi Kumulatif}
% =============================================================================
Gambar \ref{fig:dmc} dan \ref{fig:heatmap_all} menyajikan kurva akumulasi massa ganda dan matriks korelasi harian seluruh produk.

\begin{figure}[htbp]
    \centering
    \includegraphics[width=0.85\textwidth]{07_kurva_massa_ganda_harian.png}
    \caption{Kurva Massa Ganda (Double-Mass Curve) Akumulasi Curah Hujan Harian terhadap AWS IoT.}
    \label{fig:dmc}
\end{figure}

\begin{figure}[htbp]
    \centering
    \includegraphics[width=0.85\textwidth]{08_heatmap_korelasi_harian_semua_produk.png}
    \caption{Matriks Korelasi Rank Spearman Harian: 8 Produk Presipitasi Kebumen dan AWS IoT.}
    \label{fig:heatmap_all}
\end{figure}

% =============================================================================
\section{Kesimpulan}
% =============================================================================
\begin{enumerate}[leftmargin=*]
    \item \textbf{Produk Terbaik Harian}: \texttt{NASA GPM IMERG} terbukti sebagai produk satelit terbaik dalam mengestimasi curah hujan harian di Kebumen ($r = 0.619, \rho = 0.642, \text{KGE} = 0.548, \text{MAE} = 7.64\text{ mm/hari}$).
    \item \textbf{Efek Multi-Temporal}: Agregasi 24-jam meningkatkan akurasi korelasi secara signifikan ($+49.5\%$), mengonfirmasi bahwa agregasi waktu efektif mereduksi *noise* sub-harian dan ketidaksesuaian spasial.
    \item \textbf{Dinamika Diurnal Tropis}: Hujan konvektif di Kebumen mencapai intensitas maksimum pada sore hari (15:00--18:00 WIB), bertepatan dengan penurunan suhu permukaan pasca-puncak insolasi surya.
\end{enumerate}

\end{document}
